\printedpage{115}
% head: A. Grothendieck : Dérivateurs, Ch. XII, Édition des Universités de Montpellier II et Paris VII
% foot: Version 30 mars 2003

\msp{page 171}

\section*{VI\quad Criterion for $W_\infty$-asphericity of the $\Omega^n(X,x)$.}

\subsection*{1\quad Cogitations on $W_n$.}

We are going to define, for $n \in \mathbf{N}$,
\[
W_n \subseteq \mathrm{Fl}(\mathrm{Cat})
\]
by the following conditions (where $f : X \longrightarrow Y$ is an arrow in Cat).

\noindent\begin{tabular}{@{}l c l >{\RaggedRight\arraybackslash}p{0.64\textwidth}@{}}
$f \in W_0$ & $\Longleftrightarrow$ & a) & $\pi_0(f)$ bijective.\\
 & $\Longleftrightarrow$ & b) & For every presheaf of sets $F$ \textit{constant} on $Y$, $H^0(Y,F) \longrightarrow H^0(X,f^*F)$ is bijective.\\
 & $\Longleftrightarrow$ & c) & The functor $F \longmapsto f^*F$ from presheaves of sets \textit{constant} on $Y$ to constant sheaves of sets on $X$ is fully faithful.
\end{tabular}

\texttt{[The terms sheaf or presheaf are used interchangeably here.]}

\textbf{NB} If $F_0$ is a set, $F_{0\,Y}$ the constant sheaf of sets defined by $F$ on $Y$, we have
\[
\begin{aligned}
H^0(Y,F_{0\,Y}) &\simeq \HomEns(\pi_0(Y),F_0) && \text{(functorial in $Y$)}\\
\mathrm{Hom}_{Y^\wedge}(F_{0\,Y},G_{0\,Y}) &\simeq \HomEns(\pi_0(Y),\HomEns(F_0,G_0)) && \text{(functorial in $Y$),}
\end{aligned}
\]
which implies that a) $\Longrightarrow$ c) $\underset{\substack{\text{taking}\\ \mathrm{Hom}_{Y^\wedge}(e_Y,F)}}{\overset{\text{tautologically}}{\Longrightarrow}}$ b) $\Longrightarrow$ a).

\noindent\begin{tabular}{@{}l c l >{\RaggedRight\arraybackslash}p{0.64\textwidth}@{}}
$f \in W_1$ & $\Longleftrightarrow$ & a) & $\pi_0(f)$ bijective, and for every $x \in \Ob X$, $\pi_1(f,x) : \pi_1(X,x) \longrightarrow \pi_1(Y,f(x))$ is bijective. (NB It suffices to check this for \textit{one} $x$ in each connected component of $X$.)\\
 & $\Longleftrightarrow$ & a$'$) & $\Pi(X) \xrightarrow{\Pi f} \Pi(Y)$ is an equivalence of categories.\\
 & $\Longleftrightarrow$ & b) & $H^0(Y,F) \longrightarrow H^0(X,f^*F)$ bijective for every sheaf of sets \textit{locally constant} on $Y$, and $H^1(Y,F) \longrightarrow H^1(X,f^*F)$ bijective if in addition $F$ is a sheaf of groups.\\
 & $\Longleftrightarrow$ & b$'$) & As in b), but for $H^1$ we restrict ourselves to $F$ \textit{constant}, and require $H^1(f)$ \textit{surjective} only.\\
 & $\Longleftrightarrow$ & c) & The functor $F \longmapsto f^*F$ from the category of locally constant sheaves of sets on $Y$ to that of locally constant sheaves on $X$ is an equivalence of categories.
\end{tabular}

\clearpage

\printedpage{116}
% head: A. Grothendieck : Dérivateurs, Ch. XII, Édition des Universités de Montpellier II et Paris VII
% foot: Version 30 mars 2003

\msp{page 172}

\textbf{NB} The group $\pi_1(X,x)$ is defined via the theory of the fundamental groupoid. We know that it is in canonical bijection with $\pi_0(\Omega(X,x))$, where $\Omega(X,x) = \Ch(X;x,x)$ (which allows us to have a monoidal composition law on $\Omega(X,x)$, giving the group law on $\pi_0$), or at one's choice $\Omega(X,x) = \Ch_\infty(X;x,x)$ (more convenient in certain situations, but not directly giving the composition law in $\pi_0(\Omega(X,x))$).

I am going to set, for $n \in \mathbf{N}$,
\[
\pi_n(X,x) = \pi_0(\Omega^n(X,x)),
\]
$\Omega^n$ denoting the $n$-th iterate of the functor $(X,x) \longmapsto \Omega(X,x)$ from the category $\mathrm{Cat}_\bullet$ of pointed objects of Cat to itself. I prefer here to work with the \underline{Ch} theory instead of $\Ch_\infty$. Sooner or later we will have to give a comparison statement \dots

\textbf{Theorem 1} (\footnote[112]{Probably \textit{false}. We should have a) $\Longrightarrow$ b), but b) $\Longrightarrow$ a) \texttt{[is]} probably false already for $n = 2$ (except in the case $X$, $Y$ 1-connected). Cf. corrected statement \texttt{[p. 178]}.}). \textit{Let $f : X \longrightarrow Y$ be in} Cat\textit{, and let $n \in \mathbf{N}$, $n \geq 2$, $S$ a subset of} $\Ob X$ \textit{such that the composite $S \longrightarrow \Ob X \longrightarrow \pi_0(X)$ is} \texttt{[an]} \textit{epimorphism. The following conditions are equivalent.}

\begin{itemize}
\item[\textit{a)}] \textit{$\pi_i(f,x)$ an isomorphism for every $x \in \Ob X$ and every $i \in [0,n]$.}

\item[\textit{a$'$)}] \textit{As in a), but with $x \in S$.}

\item[\textit{b)}] \textit{We have}
\[
\begin{array}{rcll}
H^0(Y,F) & \xrightarrow{\ \sim\ } & H^0(X,f^*F) & \textit{for $F$ a sheaf of sets}\\
 & & & \textit{locally constant on $Y$.}\\
H^1(Y,F) & \xrightarrow{\ \sim\ } & H^1(X,f^*F) & \textit{if in addition $F$ a sheaf of groups.}\\
H^i(Y,F) & \xrightarrow{\ \sim\ } & H^i(X,f^*F) & \textit{if in addition $F$ an abelian sheaf, $i \in [2,n]$.}
\end{array}
\]
\end{itemize}

\msp{page 173}

\textbf{NB} The equivalence of a) and a$'$) should be trivial, but for this we need to have defined the $\pi_n(X,x)$, for $x$ varying in $\Ob X$, as a local system on $X$, something I have not yet done. I am going to try to do it in a draft of air. First of all, the $\Omega(X,x)$ are the fibers of a perfect $W_\omega$-fibration (\footnote[113]{\textbf{NB} $W_\omega$ is the smallest localizer of Cat (satisfying W(1, 2 bis, 3)).}), fitting into the cartesian square
\[
\begin{tikzcd}
\Ch(X) \arrow[d, "{p_X=(\sigma_X,\beta_X)}"'] & \Omega(X,-) \arrow[l] \arrow[d] \\
X \times X & X. \arrow[l, "\mathrm{diag}_X"'] \arrow[ul, dotted, "i_X"']
\end{tikzcd}
\]

\clearpage

\printedpage{117}
% head: A. Grothendieck : Dérivateurs, Ch. XII, Édition des Universités de Montpellier II et Paris VII
% foot: Version 30 mars 2003

We then know that we have a functor
\[
\begin{aligned}
\Pi X &\longrightarrow \mathrm{Hot}_{W_\omega}\\
x &\longmapsto \mathrm{hot}_{W_\omega}(\Omega(X,x)).
\end{aligned}
\tag*{$(*)$}
\]
Now we have, for $i \geq 1$,
\[
\pi_i(X,x) = \pi_{i-1}(\ \underbrace{\Omega(X,x)}_{\text{pointed object}}\ ).
\]
This gives the desired local system of the $\pi_i$, if only we had a functor a little more precise than $(*)$, namely
\[
\Pi X \longrightarrow \Hotp{W_\omega} \qquad \text{(\textit{pointed} $W_\omega$-homotopy category)}
\]
(\footnote[114]{But one must also check that the $\pi_n(X,x)$ ($n \geq 1$) are functors on $\Hotp{W_\infty}$, i.e. transform $W_\infty$-equivalences into likewise. This follows from W(4a) for $W_\infty$, which implies that $(X,x) \mapsto \Omega(X,x)$ factors through $\Hotp{W_\infty}$ (as a functor with values in $\Hotp{W_\infty}$).}). Now the arrow $i_X : X \longrightarrow \Ch X$ lifting $p_X$ gives a canonical section of $\Omega(X,-)$ over $X$ -
\msp{page 174}
precisely the one which assigns to every $x \in X$ the marked point $1_x$ of $\Omega(X;x,x)$ (``empty loop at $x$''). The question is thus this:

\begin{quote}
We have a $W$-fibration $\Omega \longrightarrow X$, with a given section $\sigma$. Does this define a canonical functor % [sic] donné -> donnée, accord avec « section »
\[
X \longrightarrow \Hotp{W},
\]
turning every arrow into \texttt{[an]} isomorphism?
\end{quote}

The natural idea is to take
\[
x \longmapsto \hotp{W}(\Omega/x)
\]
by pointing $\Omega/x$ at the marked point $\sigma(x)$ of $\Omega_x \hookrightarrow \Omega/x$. The trouble is that the transition morphisms $\Omega/x \longrightarrow \Omega/x'$ (for $x \xrightarrow{u} x'$ in $X$) are not compatible with the marked points. But we have an arrow (namely $\sigma(u)$) going
\msp{page 175}
from the image $\sigma(x)/x'$ of the marked point of $\Omega/x$ to the marked point $\sigma(x')/x'$ of $\Omega/x'$. This should suffice to define the arrow
\[
(\Omega/x,\sigma(x)/x) \longrightarrow (\Omega/x',\sigma(x')/x')
\]
as an arrow in $\Hotp{W}$, with transitivity for a composite $x \xrightarrow{u} x' \xrightarrow{v} x''$ in $X$. For the moment I will confine myself
\printedpage{118}
% head: A. Grothendieck : Dérivateurs, Ch. XII, Édition des Universités de Montpellier II et Paris VII
% foot: Version 30 mars 2003
\msp{page 176}
to an indication, leaving detailed verifications for later. If need be, I will use right now that for $i \in \mathbf{N}$, $i \geq 1$,
\[
\pi_i(f,x) : \pi_i(X,x) \longrightarrow \pi_i(Y,f(x))
\]
can be interpreted as coming from a homomorphism of local systems
\[
\utpi_i(X,-) \longrightarrow f^*(\utpi_i(Y,-))
\]
on $X$. Condition a) of Theorem 1 is interpreted by saying that this homomorphism is an \textit{isomorphism} for $i \in [1,n]$, and condition a$'$) then merely recalls that it suffices to test that this is so for \textit{one} point of each connected component of $X$.

The statement of Theorem 1 retains a meaning for $n = 1$, and in this case was ``recalled'' previously (and is $\pm$ tautological). This will allow us to proceed by induction on $n$, assuming that the theorem is proved for $n-1$. This means that we can suppose that conditions a) \textit{and} b) are satisfied by $f$ for $i \in [1,n-1]$, and confine ourselves to proving that, under these conditions, condition a) (or a$'$), for $i = n$, is equivalent to condition b) for $i = n$.

\msp{page 177}

So we simply have to prove this:

\textbf{Lemma} (\footnote[115]{Probably false.}). \textit{Suppose that $\pi_i(f,x)$} \texttt{[is an]} \textit{isomorphism for $i \in [0,n-1]$ and b) satisfied for $i \in [0,n-1]$, where $n \in \mathbf{N}$, $n \geq 2$, is fixed. Then the following conditions are equivalent:}

\begin{itemize}
\item[\textit{a)}] $\pi_n(f,x) : \pi_n(X,x) \longrightarrow \pi_n(Y,f(x))$ \texttt{[is an]} \textit{isomorphism for every $x \in \Ob X$.}

\item[\textit{b)}] $H^n(Y,F) \xrightarrow{\sim} H^n(X,f^*F)$ \texttt{[is an]} \textit{isomorphism for every locally constant abelian sheaf on $Y$.}
\end{itemize}

Let us note that $f$ induces an arrow in $\mathrm{Cat}_\bullet$,
\[
\Omega(f,x) : \Omega(X,x) \longrightarrow \Omega(Y,f(x)),
\]
whence
\[
\begin{tikzcd}[row sep=small]
\pi_i(\Omega(X,x)) \arrow[r] \arrow[d, dash, "\wr"'] & \pi_i(\Omega(Y,f(x))) \arrow[d, dash, "\wr"'] \\
\pi_{i+1}(X,x) \arrow[r] & \pi_{i+1}(Y,f(x)),
\end{tikzcd}
\]
identifying with $\pi_{i+1}(f,x)$:
\[
\pi_i(\Omega(f,x)) \simeq \pi_{i+1}(f,x).
\]


\printedpage{119}
% head: A. Grothendieck : Dérivateurs, Ch. XII, Édition des Universités de Montpellier II et Paris VII
% foot: Version 30 mars 2003

Hypothesis a) of the lemma implies that $\pi_i(\Omega(X,x)) \xrightarrow{\sim} \pi_i(\Omega(Y,y))$ for $i \in [0,n-1]$. Here we are dealing with the $\pi_i$ based at the unit points. We will probably have to prove that, because this is a morphism of monoids in Cat (satisfying certain conditions?) this remains true for the $\pi_i$ at \textit{every} point $\xi$ of $\Omega(X,n)$. By the induction hypothesis, we will thus know that $\Omega(X,x) \longrightarrow \Omega(Y,y)$ % [sic] \Omega(X,n) -> \Omega(X,x), le point de base est x
\msp{page 178}
induces an isomorphism on cohomology with locally constant coefficients, up to and including dimension $n-1$.

Finally, I have $\pm$ convinced myself that the statement of Theorem 1 is false as it stands. Here is a statement that I hope is true, and that I am going to try to prove in a ``utilitarian'' spirit, using everything I know in homotopy and cohomology (and not just what I have developed ``from scratch'' in the modelizer Cat). % [sic] j'éspère -> j'espère, orthographe

\textbf{Theorem 1 corrected.} \textit{Let us consider, for a given integer $n \geq 2$, the following conditions on an arrow $f : X \longrightarrow Y$ in} Cat \textit{(or a continuous map of topological spaces, or a morphism in $\mathbf{\Delta}^\wedge$ etc. etc.):}

\begin{itemize}
\item[$a_n$)] \textit{If $Z$ is a ``homotopy fiber'' of $f$ (NB there is essentially one per connected component of $Y$), then $Z$ is $n$-connected, i.e.}
\msp{page 179}
\textit{(e.g.) $\pi_i(Z) = 0$ for $i \in [0,n]$.}

\item[$b_n$)] \textit{$\pi_0(f)$ is bijective, and $\forall\ x \in X$, $\pi_i(f,x) : \pi_i(X,x) \longrightarrow \pi_i(Y,f(x))$ is} \texttt{[an]} \textit{isomorphism for $i \in [1,n]$. (}\textbf{NB} \textit{It suffices to say this for} \textup{one} \textit{point $x$ of each connected component of $X$.)}

\item[$c_n$)] $H^i(Y,F) \longrightarrow H^i(X,f^*F)$ \texttt{[is an]} \textit{isomorphism for $F$ a locally constant sheaf on $Y$ (with $F$ abelian if $i \geq 2$, $F$ a sheaf of groups if $i = 1$).}
\end{itemize}

\textit{We then have}
\[
a_n \ \Longrightarrow\ b_n \ \Longrightarrow \text{(\footnotemark[116])}\ c_n \ \Longrightarrow \text{(\footnotemark[117])}\ a_{n-1}
\]
\footnotetext[116]{if $X$, $Y$ 1-connected, this is really proved, otherwise, there are verifications to be made (which look sorital \dots).}%
\footnotetext[117]{\textbf{NB} Even if $X$, $Y$ 1-connected, it is not clear that $\mathrm{c}_n \Longrightarrow \mathrm{b}_n$, for one must show that $\pi_n(Z) \longrightarrow \pi_n(X)$ is zero, and one only finds that the composite with $\pi_n(X) \longrightarrow H_n(X)$ is zero. I presume that there are counterexamples.}%
\textit{(and I presume that each of these implications is strict).}

\textsc{Proof}. $\mathrm{a}_n \Longrightarrow \mathrm{b}_n$ is an immediate consequence of the homotopy exact sequence for a Hot-fibration $\bar{X} \xrightarrow{\bar{f}} Y$ which factors $X \xrightarrow{f} Y$ as $f = \bar{f} \circ i$, $i \in W_\infty$. In fact, $\mathrm{a}_n$ \textit{is equivalent} to $\mathrm{b}_n$ + the surjectivity of the $\pi_{n+1}(X,x) \longrightarrow \pi_{n+1}(Y,f(x))$, which of course does \textit{not} follow from $\mathrm{b}_n$, i.e. from the bijectivity of the $\pi_i(f,x)$ for $i \leq n$ (take some $K(\pi,n+1)$). Thus the implication is indeed strict.